%
% $Log: contents.tex,v $
% Revision 1.6  2002/10/21 15:26:44  mjm
% to ohio, revised
%
% Revision 1.5  2001/08/15 18:07:03  mjm
% APPM preprint 466
%
% Revision 1.4  1999/08/09 19:55:27  mjm
% beta release
%
% Revision 1.3  1999/04/06 19:29:00  mjm
% Matts comments
%
% Revision 1.2  1999/04/06 19:22:02  mjm
% Ellens comments
%
%

%%%%%%%%%%%%%%%%%%%%%%%%%%%%%%%%%%%%%%%%%%%%%%%%%%%%%%%%%%%%%%%%%%%%%%%%%
\handout{How to Use This Packet}

This packet contains all the materials necessary to implement the good
problems program. The items included are:
\begin{description}
	\item[The Instructor's Guide.] This guide is for the main
instructor, meaning the person in charge of organizing the course. It lists
those tasks that the 
instructor must perform to implement this scheme, indicates when they must be
performed, and estimates how much time they are expected to take.
If there is more than one instructor, then several of these tasks can be
shared. 
	\item[The Teaching Assistant's Guide.]
This guide is for the teaching assistant, who runs the recitations.
It lists those tasks that the teaching assistant must perform, indicates
when they must be performed, and estimates the time required.
If the course is taught without recitations, then these tasks
fall on the instructor.
	\item[The Student's Guide.] This guide explains the method to the
students, and tells them what is expected from them.
When this program is implemented, this guide could be one side of a sheet
of paper and the schedule of good problems for the course photocopied onto
the back. 
	\item[The Handouts.] The main material for this method is a set of
six handouts. 
Each handout is two pages long, so that it can be copied double-sided onto
a single sheet of paper. 
The handout consists of a description of a particular writing skill,
motivation for the importance of this skill,
whatever explicit rules can be given,
and examples of good and bad presentation emphasizing that skill.
 It is hoped that
the students will find these handouts useful enough that they will keep
them for future reference.

The handouts could be used in any order, but we suggest the following sequence:
\begin{enumerate}
\item Laying out the Problem.
\item Flow.
\item Mathematical Symbols.
\item Logical Connectives.
\item Graphs.
\item Introductions and Conclusions.
\end{enumerate}

\end{description}

